\medskip\noindent\textbf{Supplementary Note 11. Software, hardware and run counts.}\par\smallskip
\begingroup\small\setlength{\parskip}{3pt}
Training used scvi-tools 1.4.2, PyTorch 2.13.0 with builds 2.13.0+cu130 or 2.13.0+cu126, lightning 2.6.5 and pyro-ppl 1.9.1. GPUs were NVIDIA GeForce RTX 3080 (10 GB) and NVIDIA L40S (46 GB). Hardware type was constant within each task's sweep and replication. CPU tools used PyTorch 2.4.1+cu121, harmony-pytorch 0.1.7 and Scanorama 1.7.4.

Scoring used scib 1.1.7, R 4.5.2, kBET 0.99.6 and rpy2 3.6.7. Python was 3.11.15, scanpy 1.11.5 and scikit-learn 1.9.0. Fitting and analysis used anndata 0.12.19, numpy 2.4.6, scipy 1.17.1, pandas 2.3.3, numba 0.67.0 and pynndescent 0.6.0. Scoring used anndata 0.11.4, numpy 1.26.4, scipy 1.14.1, pandas 1.5.3, numba 0.66.0 and pynndescent 0.5.13. CPU tools also used anndata 0.11.4 and numpy 1.26.4. Some tasks were scored on two computers, with the random seed of the local inverse Simpson's index fixed. Their scores were combined only after 5 latents gave identical values in all 19 metric columns on both computers, so agreement was shown on those latents only. Supplementary Note 10 describes which scores were available when analysis rules were chosen.

The sweep comprised 3,033 fits. Input pilots added 108 posterior-sample and 180 standardised-input fits. Baselines supplied 447 of 448 planned fits. Replication comprised 855 planned fits and 40 extensions, giving 895 fits. No fit failed. The direct biology check reused 783 fits, and scaling analyses required no new fits. Supplementary Table S5 gives task-level coverage. Immune and both simulations were low-signal with the batch-conditioned decoder (Supplementary Note 5).
\par\endgroup
